% ============================================================================
% 06_ranking.tex — Section 6: The Simplicity-Preservation Ranking
% Rewrite pass (2026-07): auctions are the main domain; DA is the robustness/
% generalization domain (Appendix app:da; landing owned by the appendix
% writer). All numbers wired to the merged seed-1299 datasets:
% results/merged_ranking/{auction_cells.csv,da_cells.csv,concordance.md} and
% results/traces/{traces_summary.md,traces_subsection.tex}. E3 menu and
% clock-framing cells are four-model; tab:ranking carries pooled rho medians
% with double-bootstrap 95% CIs; fig:ranking = ranking_forest.pdf (real);
% fig:ranking-framing = figure3_intervention_comparison.png (real, 2x4);
% new subsection sec:ranking-traces (+ tab:trace-dissociation).
% NOTE for integrator: clock-framing (B3) = the intervention_proxy_breitmoser
% sealed-bid cells (SMAD 6.3-13.8%, pooled rho +0.29); the SMAD 1.4-5.2%,
% p<0.001-all-four, d-up-to-1.00 cells are the TRUE closed clock
% (ascending_clock_closed, family A), per concordance.md and the R6 audit in
% appendix_prompts. The two rungs are reported separately below.
% ============================================================================

\section{The Simplicity-Preservation Ranking}\label{sec:ranking}

This section delivers the paper's central product: an empirical ranking of the
design levers of Section~\ref{sec:framework} by how well each preserves
dominant-strategy play. Auctions are the domain throughout; one lever---the
menu description---is additionally probed in a supplementary school-choice
matching experiment (deferred acceptance) reported in
Appendix~\ref{app:da}, drawn on below where it sharpens the descriptions
result. We proceed in three passes over the
lever taxonomy. Section~\ref{sec:ranking-extensive} evaluates the
extensive-form lever (family~A): replacing the sealed-bid auction with an
ascending clock. This is the lever for which simplicity theory makes its
sharpest predictions, and it sits at the top of the ranking---but it bundles
several simplifications at once. Sections~\ref{sec:ranking-descriptions}
and~\ref{sec:ranking-scaffolds} therefore decompose the bundle with
presentation-only levers that hold the extensive form fixed: descriptions of
the mechanism (family~B) and reasoning scaffolds appended to it (family~C).
Section~\ref{sec:ranking-synthesis} assembles the cells into the ranking
itself (Table~\ref{tab:ranking} and Figure~\ref{fig:ranking}), and
Section~\ref{sec:ranking-traces} closes the section by reading the agents'
own stated reasoning against their realized bids.

Throughout, $D(L,m)$ denotes the deviation of model $m$ from
dominant-strategy play under lever $L$: the mean
bid deviation from value (bid $-$ value, in dollars), with the scaled mean
absolute deviation (SMAD; normalizer $E[b^{*}] = \$24.5$,
Section~\ref{sec:metrics}) reported in parallel. Every cell-level statistic
below---bootstrap confidence intervals (2{,}000 draws, seed 1299),
Mann--Whitney and Welch tests, Cohen's $d$---is computed from the merged
run-level dataset in the replication package
(\path{results/merged_ranking/auction_cells.csv}). Two auction baseline conventions
appear in what follows, and we fix them here once. \emph{Format-level}
comparisons (sealed bid vs.\ clock) use the dedicated SPSB cells, with
per-model mean deviations of $-0.49$ (Claude), $-1.63$ (Gemini), $-3.28$
(GPT-4o), and $-5.27$ (Gemma)---cross-model mean $-\$2.67$.
\emph{Presentation and scaffold} contrasts are tested against the pooled
axis baselines run in the same grid as each treatment---with one correction
made in this draft: a provenance audit showed that the legacy ``axis-2
baseline'' cell was itself a treated clock-framing description
(Section~\ref{sec:ranking-synthesis}), so the pool comprises the axis-1 and
axis-3 baseline cells only, with per-model means of $+0.54$ (Claude),
$-1.25$ (Gemini), $-2.94$ (GPT-4o), and $-6.53$ (Gemma)---cross-model
$-\$2.55$. The two conventions differ materially for Claude, Gemini, and
Gemma, and both are recorded cell by cell in the replication file. The comparative statistic is the preservation index
$\rho$ of Section~\ref{sec:metrics}, computed within model on
the absolute-deviation scale and pooled as the median across cells:
$\rho = 1$ is perfect preservation of dominant-strategy play, $\rho = 0$ no
effect, and $\rho < 0$ a lever that \emph{backfires}. Because $\rho$ is
computed within model, the ranking is invariant to cross-model differences in
baseline levels.

\subsection{Extensive-Form Levers: The Ascending Clock}
\label{sec:ranking-extensive}

A strategy-proof mechanism and its OSP counterpart implement the same outcome
rule---the same mapping from reports to allocations and payments---but differ
in extensive form. The non-OSP version requires \emph{contingent reasoning}:
to see that truth-telling is dominant, an agent must trace through
hypothetical scenarios about others' behavior and verify that no deviation
improves payoffs in any case. The OSP version restructures the game so that
the dominant strategy is \emph{obviously} optimal: at every information set,
the best outcome from following the strategy weakly dominates the worst
outcome from any deviation, without requiring any reasoning about what others
might do \citep{li2017obviously}. Section~\ref{sec:validation} established
that LLM agents deviate from truthful play in sealed-bid strategy-proof
formats. Here we ask how much of that deviation the extensive form itself
removes, across all four model families.

\begin{figure}[htbp]
    \centering
    \includegraphics[width=\textwidth]{fig1_osp_comparison.pdf}
    \caption{Comparing base (non-OSP) and OSP mechanisms across four model
    families. Distribution of bid
    deviations from value ($\text{bid} - \text{value}$) in second-price
    sealed-bid (gray) and ascending clock (colored) auctions. Dashed lines
    indicate means. All models deviate from truthful
    play in the non-OSP sealed-bid format; the ascending clock corrects toward
    truth-telling.}
    \label{fig:osp-comparison}
\end{figure}

\paragraph{The human-anchored APV cell.}
Clock auctions are dynamic formats in which the price rises incrementally and
bidders decide whether to stay in or drop out at each price point. Although
the sealed-bid second-price format and the ascending clock share the same
dominant strategy, human participants behave quite differently across them: a
robust pattern in the experimental literature is that bidders do much better
in clock-style auctions than in sealed-bid auctions. \citet{li2017obviously}
provides an argument for why clock auctions are easier to recognize as
strategy-proof, formalized as obvious strategy-proofness, and
\citet{breitmoser2022obviousness} provide additional evidence that clock
auctions improve truthful play, with behavior under the open clock (AC, where
dropouts are observed) slightly better than under the closed clock (AC-B).

We evaluate the clock lever in two purpose-labeled cells
(Table~\ref{tab:calibration-master}). The first is an affiliated private
values (APV) environment chosen for comparability with the human evidence of
\citet{li2017obviously} and \citet{breitmoser2022obviousness}. While
affiliation could in principle change incentives through correlation in
values, the three formats we consider are strategically equivalent: in the
second-price sealed-bid (SP), the open ascending clock (AC), and the closed
clock (AC-B), it remains a dominant strategy to bid---or drop out at---one's
value. Reconstructed human bidders deviate substantially less from truthful
bidding in the clock formats than in sealed-bid SP: SMAD is $3.54\%$ in AC and
$5.83\%$ in AC-B, compared to $9.31\%$ in SP
(Figure~\ref{fig:smad-comparison}). For GPT-4o the clock advantage is even
more pronounced. In the APV clock auctions GPT-4o achieves near-truthful play,
with SMAD falling to $0.81\%$ in AC and $0.47\%$ in AC-B (a $77\%$--$92\%$
reduction relative to the human clock baselines), whereas in sealed-bid SP it
remains at $10.73\%$---close to the human level of $9.31\%$. Breaking the
decision into a sequence of simple threshold choices (``stay iff the price is
below value'') yields an order-of-magnitude improvement over the one-shot
sealed-bid interface, even though the dominant-strategy benchmark is
identical.

The direction of the residual mistakes matters as much as their size.
Figure~\ref{fig:direction} decomposes bidding errors by direction using a
$\pm 2\%$ tolerance band around value (dropping near-equal bids and
classifying the remainder as under- vs.\ over-bids). Despite comparable
aggregate deviations in sealed-bid SP-APV ($9.31\%$ vs.\ $10.73\%$ SMAD), the
nature of the deviations differs strikingly: human bids are predominantly
\emph{above} value ($64.7\%$ overbidding vs.\ $35.3\%$ underbidding), while
GPT-4o bids are predominantly \emph{below} value ($92.4\%$ underbidding vs.\
$7.6\%$ overbidding)---humans exhibit competitive overbidding even under an
incentive-compatible rule, while GPT-4o is overly conservative in the
sealed-bid interface. The clock aligns not only the magnitude but the
\emph{shape} of behavior: in AC, humans and GPT-4o display an almost perfectly
symmetric under/over split ($49.2/50.8$ vs.\ $50/50$), with GPT-4o
concentrating far more mass near truthful bidding ($71.7\%$ within $\pm 2\%$
vs.\ $26.8\%$ for humans). We return to what this direction reversal does---%
and does not---imply for the ranking in Section~\ref{sec:humans}.

\paragraph{The four-model IPV cell.}
The second cell runs the same lever in the independent private values (IPV)
environment of Table~\ref{tab:calibration-master}, across all four model
families; its purpose is identification rather than human comparability.
In the SPSB baseline (Figure~\ref{fig:osp-comparison}, left column, gray), all
four models deviate from truthful bidding, with mean deviations of
$\mu = -0.49$ (Claude), $-1.63$ (Gemini), $-3.28$ (GPT-4o), and $-5.27$
(Gemma), for a cross-model mean of $-\$2.67$; these are the dedicated SPSB
cells of the baseline declaration above. The negative sign indicates
systematic \emph{underbidding}---the mirror image of the human overbidding
documented in Section~\ref{sec:validation}. The distributions also exhibit
long tails: some bids deviate by \$10--20 from value, cases in which the
model substantially misplays the auction.

The closed ascending clock corrects deviations toward truthful play in every
model. Mean deviations compress to $\mu = -0.05$ (Claude), $-1.24$ (Gemini),
$-0.45$ (GPT-4o), and $-0.33$ (Gemma)---a cross-model mean of roughly
$-\$0.5$---and SMAD falls to $1.4\%$, $5.2\%$, $1.9\%$, and $1.9\%$ against
pooled-axis baselines of $9.5\%$, $11.7\%$, $11.4\%$, and $20.7\%$. The shift
is significant for every model (Welch $p < 0.001$ in all four cells), with
effect sizes up to Cohen's $d = 1.00$ (Gemma). The improvement is most
dramatic for the models with the worst SPSB performance: Gemma's mean
deviation falls from $-5.27$ to $-0.33$, a near-complete correction. Beyond
the mean shifts, the clock \emph{eliminates the long tails} visible in the
SPSB distributions: the catastrophic deviations effectively disappear. This
variance reduction may be as practically important as the mean
correction---OSP does not merely improve average play, it also removes the
worst-case failures. Because this cell is run under IPV, the improvement
cannot be informational.\footnote{In a common- or affiliated-values setting,
the clock format could improve play by revealing information about others'
values through their dropout decisions. Under IPV, no such information
exists, so the improvement must be purely cognitive.}
Two provenance notes complete this cell's accounting. The environment audit
(plan E2) found that the legacy \path{ascending_clock_closed} runs pair an
affiliated-values \emph{description} with independent private-value
\emph{draws} (the configuration's value model sets the common component to
zero), so the informational channel was already absent in the data---but the
prompt and the environment disagreed. A dedicated clean IPV cell
(\path{ascending_clock_ipv_closed}: IPV description, IPV draws, $\$0.5$
tick, seed 1299, $K = 50$) now exists for GPT-4o and Gemma and reproduces
the result exactly---SMAD $1.5\%$ (GPT-4o) and $1.8\%$ (Gemma) against
sealed baselines of $12$--$27\%$---so the identification claim does not rest
on the mislabeled cells. The remaining IPV-clock cells for Claude 3.5 Haiku and Gemini
2.0 Flash await native-API access (both model snapshots have been retired
from the aggregator used for this draft's re-runs); given the exact
replication in the two available families and the direction of every legacy
cell, we do not expect the remaining two cells to change the rung.

\paragraph{The bridge: which margin does the work?}
In summary, LLMs deviate from dominant-strategy play in the sealed-bid
strategy-proof auction, and the OSP ascending clock substantially corrects these
deviations. As a design lesson, however, the extensive-form result is coarse,
because OSP changes several things at once: it changes the game tree, it
coarsens the action space to simple binary or menu choices, and it makes the
safety property of truth-telling salient at every decision point. The clock
eliminates the need for contingent reasoning, reduces each decision to a
stay/exit choice, and makes the connection between action and outcome
immediate. Which of these simplifications actually drives the improvement?
The theory does not need to answer this question---any OSP implementation
bundles them---but a designer does, because the margins have very different
costs: re-engineering the extensive form of a deployed mechanism is
expensive, while changing its description, or the reasoning it invites, is
nearly free. The next two subsections decompose the bundle with
presentation-only levers that hold the sealed-bid extensive form fixed:
descriptions of the mechanism (family~B,
Section~\ref{sec:ranking-descriptions}) and reasoning scaffolds appended to
it (family~C, Section~\ref{sec:ranking-scaffolds}).

\subsection{Descriptions: Restating the Rules versus Exposing the Safety
Invariant}\label{sec:ranking-descriptions}

Family~B holds the mechanism fixed and varies only how it is described.
Recent studies ask whether behavior in strategy-proof mechanisms depends on
how the mechanism is \emph{presented}, holding the underlying incentives
fixed. \citet{gonczarowski2023strategyproofness} propose
strategyproofness-exposing descriptions, including an extensive-form
\emph{menu} description that compresses the dominant-strategy logic of a
second-price auction into a short menu statement intended to make
strategy-proofness salient; in their laboratory experiments, however, this
menu framing does not meaningfully change human bidding relative to the
traditional SPSB description, and \citet{gonczarowski2024describing} likewise
find that menu-based explanations of DA improve measured \emph{understanding}
while average behavioral effects remain modest. In contrast,
\citet{breitmoser2022obviousness} show that \emph{reframing} SPSB as a
dynamic ascending clock substantially increases truthful play by human
subjects. The taxonomy of Section~\ref{sec:framework} accordingly splits the
family into three sub-levers---menu restatement (B1), safety-exposing
descriptions (B2), and clock-framing (B3)---and they turn out to behave very
differently.

\paragraph{Menu restatement (B1): null on average, sign-mixed underneath.}
The menu restatement has no consistent effect on bidding. Across the four
models, its effects are small and sign-mixed relative to each model's pooled
axis baseline (Figure~\ref{fig:ranking-framing}, top row): Gemma is a clean
null (mean deviation $-4.50$ vs.\ baseline $-4.54$; Welch $p = 0.94$), Claude
moves slightly \emph{away} from truthful play ($+0.48$ to $+1.37$,
$\Delta = +0.88$ above value, $p = 0.042$, $d = 0.24$), GPT-4o moves
\emph{away} as well ($-2.78$ to $-4.10$, $\Delta = -1.32$, $p = 0.003$,
$d = -0.37$), and Gemini alone moves \emph{toward} truthful bidding ($-2.88$
to $-1.13$, $\Delta = +1.75$, $p < 0.001$, $d = 0.50$). Pooling across cells,
the restatement is a null on average: a two-sided sign test of the B1 cells
against their baselines cannot reject no effect ($p = 0.73$;
Section~\ref{sec:ranking-synthesis}). The pattern also survives payment-rule
variation: in GPT-4o's first- and third-price variants of the same grid, the
menu cells are indistinguishable from their baselines ($p = 0.63$ and
$0.51$). We therefore read the menu cell as consistent with the human
null \emph{on average} \citep{gonczarowski2023strategyproofness,
gonczarowski2024describing}, with the caveat that the average conceals
genuine model-level heterogeneity---a point the human experiments, run on a
single population, cannot speak to. The menu family turns out to be more
interesting in the supplementary matching experiment, where it splits into two
sub-variants with opposite effects; we return to this in
Section~\ref{sec:ranking-synthesis}.

\paragraph{Clock-framing (B3): describing the bid as a clock exit.}
The clock-framing description keeps the one-shot sealed-bid elicitation---the
agent still submits a single number---but describes that number as the
threshold at which an automatic proxy would exit an ascending clock on the
bidder's behalf. This is a description, not a mechanism change: no prices are
posted, no dropouts occur, and no information is revealed; it should not be
conflated with the AC/AC-B formats of Section~\ref{sec:ranking-extensive}.
The two rungs are empirically distinct cells in the merged dataset: the
clock-framing cells record exactly one sealed bid per bidder per auction,
while the true clock cells record per-tick stay/exit decisions
(Appendix~\ref{app:prompts}). Against the corrected pooled baselines
(Section~\ref{sec:ranking-synthesis}), clock-framing moves play
significantly in every model family but not always in the same direction
(Figure~\ref{fig:ranking-framing}, bottom row): SMAD falls from $12.1\%$ to
$6.3\%$ for GPT-4o and from $26.8\%$ to $13.8\%$ for Gemma---each
roughly halving its deviation---and from $8.3\%$ to $7.2\%$ for Claude (all
Welch $p < 0.001$), while Gemini moves the \emph{wrong} way, from $5.1\%$ to
$9.5\%$ ($p < 0.001$). The sign agrees with the human evidence that clock
reframing increases truthful play \citep{breitmoser2022obviousness} in three
of the four families, and the exception is coherent rather than noisy:
Gemini is also the one model that the two-stage clock-exit variant below
sharply harms. Pooled, the rung is modest and dispersed---median
$\rho = +0.30$ ($95\%$ CI $[-0.77, +0.52]$, the width driven by the Gemini
reversal)---recovering, where it works, a third to a half of the gap to
truthful play, where the true extensive-form clock recovers nearly all of it
(SMAD $1.4$--$5.2\%$; Section~\ref{sec:ranking-extensive}). Words that
borrow the clock's framing can help, and for some models help substantially;
the clock itself helps every model, and helps far more.

The rung comes with a measured warning about paraphrase sensitivity. A
provenance audit of the legacy run grid (documented in the replication
package) uncovered a \emph{second} clock-framing text that had been
mislabeled as an axis-2 baseline cell: a two-stage description that presents
the sealed bid as an ``exit price'' at which an ascending price clock would
remove the bidder. Its manipulation check passes---clock-and-exit language
appears in $14.3\%$ of treated traces against $0\%$ in the corrected
baselines (wild-cluster $p = 0.002$)---but its behavioral effects, unlike
the canonical clock-framing text above, are violently heterogeneous: it
nearly fixes Gemma (mean deviation $-6.53$ to $-0.59$), sharply
\emph{backfires} for Gemini ($-1.25$ to $-6.12$), and leaves Claude and
GPT-4o unmoved, with GPT-4o inert in language as well as bids (zero echo of
the framing), so that the pooled effect is a null that conceals offsetting
model-level responses ($p = 0.88$ pooled). The two texts agree on one
regularity---both harm Gemini, the family most reactive to clock
vocabulary---but otherwise produce different model-level profiles from the
same clock analogy: the canonical framing helps three families, the
two-stage variant helps only Gemma. This is direct evidence that
description effects are sensitive to wording in a way that extensive-form
changes are not (the true clock helps every model; Section~%
\ref{sec:ranking-extensive}), and a caution against reading any single
description cell, including ours, as \emph{the} effect of its family.

\paragraph{Safety-exposing descriptions (B2): stating the invariant.}
The safety-exposing sub-family does not restate the rules at all; it states
the structural property of the mechanism that makes truth-telling safe. The
key step in \citet{vickrey1961counterspeculation}'s original argument for the
strategy-proofness of the second-price auction concerns being
\emph{pivotal}: your bid determines only \emph{that} you win, not \emph{what}
you pay. This inspires our \emph{Payoff Safety} intervention for the auction
setting: ``Think of your bid as setting your maximum price. Your bid
determines IF you win, not WHAT you pay. If you win, you pay what others bid,
not what you bid.'' Its analogue in the supplementary matching experiment, \emph{Rejection
Safety}, states the property that made deferred acceptance preferable to the
Boston mechanism in practice \citep{abdulkadirouglu2003school}: ``Rejections
only redirect, never eliminate. Your ranking determines the \emph{order} in
which you are considered, not \emph{whether} you are considered.'' Neither
intervention scaffolds a particular mode of reasoning: they do not ask the
agent to trace through contingencies, simulate the algorithm, or reason about
opponents. They simply state the key property of the mechanism that makes it
strategy-proof.

In auctions, \emph{Payoff Safety} reduces the mean deviation from
$\mu = -2.67$ to $-1.53$ (Figure~\ref{fig:descriptions}), with the deviation
shrinking in magnitude for every model: Claude from $-0.49$ to $+0.30$,
Gemini from $-1.63$ to $-0.38$, GPT-4o from $-3.28$ to $-1.88$, and Gemma
from $-5.27$ to $-4.17$.\footnote{The Claude endpoint is $+0.30$
(overbidding) in the raw logs; the ES manuscript reported it as $-0.30$. The
magnitude is identical and the qualitative claim---Claude ends closest to
truthful---is unchanged, but the sign correction means Claude's residual
error flips from under- to over-bidding under this description.} Against the
axis-specific baselines the improvement is significant for Gemini
($p < 10^{-30}$, $d = 1.85$) and GPT-4o (Mann--Whitney $p < 10^{-5}$); for
Claude the mean effect is null while dispersion halves (SMAD $12.0\%$ to
$4.8\%$); Gemma is the one negative mark, moving away from its axis-2
baseline---though that baseline is anomalously low relative to Gemma's other
baselines, a caveat we quantify in Section~\ref{sec:ranking-synthesis}. For
comparison, the ascending clock achieves $\mu \approx -0.5$
(Section~\ref{sec:ranking-extensive}); the description closes roughly half
the gap.

That a brief description of a mechanism's incentive properties can nearly
replicate OSP's effect---without changing the extensive form---is a
surprising result. It suggests that the cognitive barrier for LLMs in non-OSP
mechanisms is not in \emph{computing} the dominant strategy given the game
tree, but in understanding the mechanism well enough to \emph{see that
truth-telling is safe}. OSP mechanisms resolve both problems simultaneously:
they simplify the game tree and make the safety guarantee transparent. Our
results suggest it is the latter that does most of the work, at least for
LLMs.

Descriptions, then, are two levers, not one---and the split is not between
``more text'' and ``less text'' but between what the text \emph{asserts}.
Restating the rules, even in the theoretically grounded menu form, is null on
average and sign-mixed underneath; exposing the safety invariant recovers
most of what full OSP implementation delivers, in every model whose baseline
leaves room to improve. The menu cell still functions as a placebo in the
aggregate---not every additional sentence about the mechanism moves behavior,
so the large effects of B2 cannot be attributed to generic prompt
sensitivity---but it is a noisier placebo than a single-population experiment
would suggest. What the effective descriptions share, and the restatement
lacks, is that they \emph{assert} the invariant that makes truth-telling safe
rather than leaving the agent to derive it. The supplementary matching
experiment sharpens this reading to a within-family contrast, discussed in
Section~\ref{sec:ranking-synthesis}: a menu description that adds the
invariance statement behaves like a safety description, while a menu
description that merely restates the computation backfires.

\begin{figure}[htbp]
    \centering
    \includegraphics[width=\textwidth]{fig5_mechanism_description.pdf}
    \caption{Safety-exposing description (B2): \emph{Payoff Safety} in
    second-price auctions (``your bid determines IF you win, not WHAT you
    pay''), across the four model families. Baseline in
    gray, intervention in color; dashed lines indicate means. The intervention
    improves play in every model, closing roughly half the gap to the full OSP
    result of Figure~\ref{fig:osp-comparison}.}
    \label{fig:descriptions}
\end{figure}

\begin{figure}[htbp]
    \centering
    \includegraphics[width=0.9\linewidth]{figure3_intervention_comparison.png}
    \caption{Presentation-only descriptions of the same SPSB mechanism,
    across all four model families ($2 \times 4$ panels; deviations in
    dollars). \emph{Top row:} menu restatement (B1) vs.\ the pooled axis
    baseline---small, sign-mixed shifts (Gemma null, Claude and GPT-4o
    slightly away from truthful play, Gemini toward it). \emph{Bottom row:}
    clock-framing (B3)---a same-signed improvement in every model,
    significant in three of four. Only LLM simulation data are shown.}
    \label{fig:ranking-framing}
\end{figure}

\subsection{Reasoning Scaffolds: Contingencies Help, Strategizing Hurts}
\label{sec:ranking-scaffolds}

Family~C holds both the mechanism and its rule description fixed and appends
a reasoning scaffold. Following \citet{li2024designing}, who decomposes the
cognitive demands of mechanisms into three dimensions, we scaffold thinking
\emph{contingently} about unobserved moves by other players (C1),
\emph{planning} for one's own future moves (C2), and reasoning about other
players' \emph{beliefs} (C3). Each scaffold directs reasoning without
revealing the dominant strategy; each modifies only the rule-explanation
block of the prompt while the mechanism, value distributions, and response
format are unchanged (full prompt texts in Appendix~\ref{app:prompts}).
Prospect-theory-inspired preference framings (family~D) are reported in
Appendix~\ref{app:prospect}; they form the ranking's bottom rungs.

\subsubsection{Contingent reasoning (C1): the payoff tree}
\label{sec:ranking-contingent}

\citet{li2024designing} identifies contingent reasoning as the core cognitive
demand of non-OSP mechanisms. To see that truth-telling is dominant in SPSB,
a bidder must consider each possible profile of opponent bids and check that
no deviation improves payoffs in any case---a ``case-by-case comparison,
calculating payoffs for each profile of opponent bids''
\citep{li2024designing}. The ascending clock eliminates this demand: at each
price, the worst case from following the dominant strategy weakly dominates
the best case from any deviation, without reference to opponents. The
question is whether scaffolding contingent reasoning---without changing the
extensive form---can partially recover these gains. Our \emph{Payoff Tree}
intervention presents the agent with the explicit decision tree of outcomes:
``PATH A: your bid exceeds others $\to$ you WIN, and pay the second-highest
bid. PATH B: your bid falls below $\to$ you LOSE, and pay nothing. Key
insight: your bid determines which path you are on, not the payment in Path
A.''

\begin{figure}[htbp]
    \centering
    \includegraphics[width=\textwidth]{fig2_contingent.pdf}
    \caption{Contingent reasoning: \emph{Payoff Tree} intervention (C1).
    Second-price auction bid deviations
    ($\text{bid}-\text{value}$) across the four model families; baseline SPSB
    in gray, intervention in color. The payoff-tree scaffold consistently
    improves play across models, roughly halving the distance from truthful.}
    \label{fig:contingent}
\end{figure}

In auctions, \emph{Payoff Tree} reduces the mean deviation from $\mu = -2.67$
to $\mu = -1.13$, roughly halving the distance from truthful play
(Figure~\ref{fig:contingent}). The improvement is consistent across models:
GPT-4o moves from $-3.28$ to $-0.86$, Gemma from $-5.27$ to $-2.34$, and
Gemini from $-1.63$ to $-0.53$; Claude, already close to truthful ($-0.49$),
is the one model showing null improvement. The scaffold closes roughly half
the gap to the OSP benchmark (clock $\mu \approx -0.5$) without changing the
extensive form.

Laying out the payoff tree does not tell models \emph{what} to do; it tells
them \emph{what happens} under each action. This is precisely the cognitive
gap between OSP and non-OSP mechanisms: in OSP mechanisms the extensive form
makes the payoff-relevant contingencies visible at each decision point,
whereas in non-OSP mechanisms the agent must construct them herself. The
\emph{Payoff Tree} scaffold does this construction for the agent, and the
consistent improvement across models suggests that the
difficulty of constructing payoff-relevant contingencies is a real constraint
that binds on LLM agents.

\subsubsection{Belief formation (C3): making opponents salient}
\label{sec:ranking-beliefs}

The third axis is belief formation. \citet{borgers2019strategically}
introduce \emph{strategic simplicity}: a mechanism is strategically simple if
optimal play does not depend on higher-order beliefs about what opponents
believe. Strategy-proof mechanisms are trivially strategically simple---the
dominant strategy is independent of beliefs entirely. Scaffolding belief
formation should therefore be irrelevant at best. But if models are already
treating these mechanisms as Bayesian games, attempting to best-respond to
perceived competition rather than playing the dominant strategy, then making
beliefs more salient could amplify mistakes. We test two levels of belief
scaffolding: \emph{First-Order Beliefs} (``What do you expect
the other players to do? \ldots Does this affect what you should do?'') and
\emph{Second-Order Beliefs} (``What do the other players think YOU will do?
\ldots And does that, in turn, affect what you should do?'').

\begin{figure}[htbp]
    \centering
    \includegraphics[width=\textwidth,height=0.82\textheight,keepaspectratio]{fig4_beliefs.pdf}
    \caption{Belief-formation interventions (C3): second-price auction bid
    deviations under first-order and second-order belief scaffolds, across the
    four model families. Baseline distributions in gray, interventions in
    color. Belief scaffolds almost uniformly worsen play, with second-order
    beliefs causing the most damage. Gemini is the notable exception, showing
    near-immunity to belief scaffolding.}
    \label{fig:beliefs}
\end{figure}

Belief scaffolding almost uniformly worsens play
(Figure~\ref{fig:beliefs}). In auctions, \emph{First-Order Beliefs} moves the
cross-model mean deviation from $\mu = -2.67$ to $-3.04$ and
\emph{Second-Order Beliefs} to $-3.47$: models prompted to think about
opponents shade further below value, undercutting perceived competition.
Interestingly, Gemini is almost completely immune to belief scaffolding, while
Gemma is the most susceptible to this failure mode. These results suggest that
LLM reasoning styles are not
monolithic; the heterogeneity across model families in susceptibility to
belief-based reasoning may itself be relevant to mechanism design when
markets serve diverse artificial agents.

\subsection{The Ranking Assembled}\label{sec:ranking-synthesis}

Table~\ref{tab:ranking} assembles the preceding cells into the paper's
central object. For each lever we report the raw auction deviations and
the preservation index $\rho$ (Section~\ref{sec:metrics}), computed within
model and pooled as the median across the available cells,
with $95\%$ confidence intervals from a double bootstrap that resamples cells
and propagates each cell's own bootstrap uncertainty (2{,}000 draws, seed
1299). Because $\rho$ is computed against each model's own baseline, the
ranking is invariant to disputes about cross-model levels; it asks only how
much of each model's own deviation a lever removes.

\paragraph{The tier structure, tested.}
Our pre-committed reading of the data is a \emph{tier structure}---a partial
order, not a strict total order:
\begin{equation*}
\underbrace{\text{A}}_{\substack{\text{extensive form:}\\ \text{ascending clock}}}
\;\succ\;
\underbrace{\{\text{B2},\,\text{B3},\,\text{C1}\}}_{\substack{\text{safety descriptions,}\\ \text{clock-framing, payoff tree}}}
\;\succ\;
\underbrace{\text{baseline} \,\approx\, \text{B1}}_{\text{menu restatement}}
\;\succ\;
\underbrace{\text{C3}}_{\substack{\text{belief}\\ \text{scaffolds}}}
\;\succ\;
\underbrace{\text{D-averse.}}_{\substack{\text{risk-averse}\\ \text{persona}}}
\end{equation*}
The concordance statistics now exist, and they support most---not all---of
this structure. Treating each model as a judge ranking the levers, Kendall's
$W$ is $0.70$ in auctions ($p = 0.004$): the lever ordering is far from
independent across models. (This statistic is computed on the corrected
baselines described below; correcting the contaminated axis-2 ``baseline''
cell \emph{raised} auction concordance from $0.43$ to $0.70$---the mislabeled
treatment had been injecting noise into every axis-2 contrast.) Exact
binomial sign tests on adjacent tiers confirm three of the four boundaries:
the extensive form beats the strong presentation tier in $7/8$ cells
($p = 0.070$), the strong presentation tier beats the restatement tier in
$8/8$ ($p = 0.008$), and the strategizing scaffolds beat the risk-averse
persona in $8/8$ ($p = 0.008$).
The boundary that \emph{fails} is baseline-vs-menu against the strategizing
scaffolds: restatements beat lookahead and beliefs in only $4/8$ cells
($p = 1.0$). The failure is informative rather than noise: the menu
restatement is not a clean placebo tier but a sign-mixed family, null on
average (two-sided placebo sign test against baseline: $3/8$ cells above
zero, $p = 0.73$) yet with genuine model-level backfires (GPT-4o's auction
menu cell, Welch $p = 0.003$) that land it below some strategizing cells.

\paragraph{The sharpened story: invariance content is the active ingredient.}
The matching-domain menu cells (new in this draft; Appendix~\ref{app:da})
resolve what the sign-mixed auction menu leaves ambiguous, by splitting the
restatement family in two. A menu description that \emph{includes the
invariance statement}---``your ranking cannot change your obtainable set,
only which school you get within it''---behaves like a safety-exposing
description: pooled error falls from $4.2\%$ to $1.8\%$, with significant
improvements for the two weak-baseline models (Claude $5.8\%$ to $2.1\%$,
Mann--Whitney $p = 10^{-4}$; Gemini $7.6\%$ to $0.3\%$, $p < 10^{-4}$) and
nulls for the strong ones. A menu description that restates the
\emph{computation} without the invariance statement \emph{worsens} play
(pooled $4.2\%$ to $6.9\%$; GPT-4o $1.0\%$ to $4.2\%$, Gemma $2.6\%$ to
$13.7\%$). Together with the auction results, the pattern across all of
family~B is consistent: safety-exposing text (B2, pooled $\rho = +0.90$) and
the invariance-carrying menu variant ($+0.65$) help; mechanical restatement
(auction menu $-0.22$, mechanics-only DA menu $-1.67$) is null-to-harmful.
The active ingredient of an effective description is not restating the
mechanism but asserting the invariant that makes truth-telling safe. The
human experimental record, read through the same split, independently
supports it: invariance-carrying descriptions improve human understanding
and (modestly) behavior in matching mechanisms
\citep{gonczarowski2024describing,katuscak2024simpler}, while
mechanics-only descriptions of a strategy-proof mechanism have been found
to \emph{reduce} truth-telling in the field \citep{guillen2018effectiveness}
---the same backfire our mechanics-only cell produces.

\paragraph{Within-tier order and an honest caveat.}
We deliberately do not order the safety description and the payoff tree within
the strong tier: their pooled preservation indices ($\rho = +0.78$ and
$+0.59$) have overlapping confidence intervals, and the payoff tree beats the
safety description in only $3/4$ of the auction cells, so the data do not
support a strict order between them. Clock-framing, now estimated on
all four models, is a reliable but modest member of the tier
($\rho = +0.29$ $[+0.17, +0.45]$). One caveat from the concordance
analysis qualifies individual cells without moving any rung: what
previously appeared as an anomalously good Gemma axis-2 baseline (SMAD
$8.6\%$ vs.\ $20.7\%$ pooled across its axes) turns out to have a mechanical
explanation: a provenance audit of the run templates shows the legacy
``axis-2 baseline'' cell was not a plain SPSB description at all but a
two-stage sealed-bid-as-clock-exit description---a clock-framing variant
that Gemma responds to strongly (Section~\ref{sec:ranking-descriptions}).
The pooled axis baseline used throughout therefore excludes that cell
(axis-1 and axis-3 cells only), the axis-2 treatment contrasts are computed
against the corrected pool, and the cell itself enters the description
family as a treatment.
\paragraph{The frontier: the ranking compresses, and the failures change kind.}
The full lever grid now also exists for four frontier models (gpt-5-mini,
gpt-5, Claude Sonnet 5, Gemini 2.5 Flash; Appendix~\ref{app:ablations}), and
it turns the frontier scope condition from an extrapolation into a
measurement. Three facts summarize it. First, \emph{the ranking's positive
rungs have nothing left to do}: all four models bid essentially exactly
truthfully at every sealed-bid baseline (mean deviations $0.00$--$0.30$,
SMAD $0$--$2.3\%$), so the safety descriptions, payoff tree, and belief
scaffolds leave play unchanged---the preservation index is undefined as
$D(\text{baseline}) \to 0$, exactly the compression anticipated in
Section~\ref{sec:discussion}. Second, \emph{the backfire rungs retain their
teeth, heterogeneously}: the menu restatement---inert to slightly harmful in
the incumbent grid---collapses Claude Sonnet 5 from exact truthfulness to a
mean deviation of $-\$8.71$ (SMAD $35.5\%$, Cohen's $d = -3.8$,
$p < 0.001$) and degrades Gemini 2.5 Flash (SMAD $2.3\% \to 9.1\%$,
$p < 0.001$), while leaving gpt-5 and gpt-5-mini exactly truthful. The
traces say exactly what broke: under the menu wording, Claude Sonnet 5
announces that ``the standard first-price auction equilibrium bid is
$(n-1)/n$ of my value'' and executes that formula flawlessly---its
$-\$8.71$ is almost exactly the $-\tfrac{1}{3}\,\mathbb{E}[v] \approx
-\$8.17$ the first-price rule implies---despite the prompt stating that
the winner pays the highest \emph{rival} bid. The unfamiliar wording did
not degrade its play into noise; it re-classified the game and retrieved
the wrong textbook solution, executed perfectly. A description with no
incentive content can still break an otherwise perfect strategic reasoner,
and which reasoner it breaks is not predictable from baseline capability. Third, what first appears as an
\emph{inversion of the top rung}---the ascending clock, which improves
every incumbent model, makes Claude Sonnet 5 \emph{worse} than its own
sealed baseline (SMAD $0.27\%$ sealed vs.\ $8.6\%$ clock, $p<0.001$, with
$75\%$ of its exits more than a dollar early)---turns out on inspection to
be the same retrieval failure as the menu cell, not a property of the
extensive form. The clock cells inherit the legacy affiliated-values
\emph{description} (Section~\ref{sec:ranking-extensive}), and Claude
Sonnet 5's exit rationales explicitly invoke the winner's curse (``others'
persistence signals possibly similar high values, risking overpayment''),
a script that is correct in common-value clocks and irrelevant here, where
the bidder knows its own value. Rerunning the same model on the clean
IPV-described clock restores near-perfect play: SMAD $2.0\%$, early exits
$3\%$, and marginal-reasoning traces (``price equals my value; drop'').
One sentence of affiliation language costs an otherwise exactly-truthful
model eight SMAD points by activating the wrong playbook---while the same
sentence leaves the incumbent models' clock gains intact, because they have
no winner's-curse script to misfire. (gpt-5 and gpt-5-mini shift by
$+0.3$--$0.6$ on the clock, mostly tick granularity, with clean marginal
reasoning in their traces.) The design lesson for mixed fleets is asymmetric:
simplifying levers are (at worst) wasted on frontier agents, but
presentation choices can still harm them, so the practitioner guidance of
Section~\ref{sec:discussion}---expose invariants, avoid gratuitous
restatements---remains binding even where the bounded-rationality regime
has closed. (The gpt-5 column carries reduced cell sizes, $K = 5$--$50$:
the model frequently omits the required response tags, and unparseable
repetitions are dropped rather than imputed; Appendix~\ref{app:ablations}
reports per-cell $n$.)

\begin{table}[htbp]
\centering
\footnotesize
\setlength{\tabcolsep}{4pt}
\resizebox{\textwidth}{!}{%
\begin{tabular}{llll cc l l}
\toprule
& & & & \multicolumn{2}{c}{Deviation $D$} & & \\
\cmidrule(lr){5-6}
Tier & Lever & Family & Theory tag & Auct.\ (\$) & DA (\%) & Pooled $\rho$ [95\% CI] & Human anchor \\
\midrule
1 & Extensive form (AC-B clock / iterative DA) & A & OSP \citep{li2017obviously,ashlagi2018stable} & $-0.5$ & 0.0 & $+0.97$ [$+0.83, +1.00$] & AC $+0.62$ [$+0.39,+0.79$]; AC-B $+0.37$ [$-0.12,+0.72$] \\
\midrule
2 & Safety description & B2 & SP-exposing \citep{gonczarowski2023strategyproofness} & $-1.53$ & 0.2 & $+0.78$ [$+0.39, +1.00$] & --- (lab prediction) \\
3$^{*}$ & Menu, invariance property (DA) & B1 & menu + invariant & --- & 1.8 & $+0.65$ [$-0.52, +0.95$] & $\surd$ agrees (sign) \citep{gonczarowski2024describing,katuscak2024simpler} \\
2 & Payoff tree & C1 & contingent reasoning \citep{li2024designing} & $-1.13$ & 1.7 & $+0.59$ [$+0.46, +0.68$] & --- (lab prediction) \\
2 & Clock-framing & B3 & OSP description \citep{breitmoser2022obviousness} & $-1.93^{\dagger}$ & --- & $+0.30$ [$-0.77, +0.52$] & $\surd$ agrees (sign, 3/4 models) \\
\midrule
3 & Baseline (SPSB / direct DA) & --- & --- & $-2.67$ & 4.2 & 0 & $\surd$ agrees \\
4 & Forward planning (lookahead) & C2 & one-step simplicity \citep{pycia2023theory} & $-1.65^{\dagger\S}$ & 7.8 / 5.1 & $-0.00$ [$-2.12, +0.26$] & --- (lab prediction) \\
3 & Menu restatement (auction) & B1 & menu description \citep{gonczarowski2023strategyproofness} & $-2.09^{\dagger}$ & --- & $-0.14$ [$-1.00, +0.27$] & $\surd$ agrees on avg.\ ($-0.07$ [$-0.12,-0.02$]) \\
4 & Second-order beliefs & C3 & strategic simplicity \citep{borgers2019strategically} & $-3.47$ & 9.1 & $-0.35$ [$-1.36, +0.03$] & --- (lab prediction) \\
4 & First-order beliefs & C3 & strategic simplicity \citep{borgers2019strategically} & $-3.04$ & 5.1 & $-0.41$ [$-0.88, -0.04$] & --- (lab prediction) \\
3$^{*}$ & Menu, mechanics only (DA) & B1 & menu, computation restated & --- & 6.9 & $-1.67$ [$-5.01, +0.57$] & $\surd$ agrees (backfire) \citep{guillen2018effectiveness} \\
\midrule
5 & Risk-averse persona & D & preference framing \citep{kahneman1979prospect} & App.~\ref{app:prospect} & App.~\ref{app:prospect} & $-1.82$ [$-5.73, -0.72$] & --- (lab prediction) \\
\bottomrule
\end{tabular}}
\caption{The simplicity-preservation ranking, rows sorted by pooled $\rho$.
$D$ = mean deviation from dominant-strategy play: bid $-$ value in dollars
for auctions (cross-model mean over GPT-4o, Claude 3.5 Haiku, Gemini 2.0
Flash, Gemma 27B), normalized Kendall-$\tau$ error in percent for the
supplementary matching experiment (Appendix~\ref{app:da}; lookahead entries are
one-/two-step). $\rho = 1 - D(\text{lever})/D(\text{baseline})$ on the
absolute-deviation scale, computed within model $\times$ domain against the
baseline conventions of Section~\ref{sec:ranking}; pooled value = median over
available cells, with double-bootstrap $95\%$ CIs (2{,}000 draws, seed 1299).
Tier = pre-committed partial order; within-tier order is not asserted.
$^{*}$New matching-domain menu cells assigned ex ante to the restatement
family (tier 3); their measured positions straddle the ranking---the
invariance-property variant performs like a tier-2 safety description, the
mechanics-only variant like a tier-4 backfire (see text).
$^{\dagger}$Intervention-grid cells tested against pooled-axis baselines
(corrected pool: axis-1 and axis-3 cells only, cross-model mean $-2.55$; see the axis-2 provenance discussion above) rather than the dedicated SPSB cells ($-2.67$).
$^{\S}$The auction C2 cell is a backward-induction \emph{framing} of the
static sealed bid (SPSB has no true planning axis); it enters the pooled
$\rho$ alongside the DA lookahead cells. Human anchor: $\surd$ = a human
experiment exists for this rung and its effect agrees in sign with the LLM
effect (reconstructed human $\rho$ [CI] shown where computable); --- = no
human experiment exists and the rung is issued as a lab prediction
(Section~\ref{sec:humans}).}
\label{tab:ranking}
\end{table}

\begin{figure}[htbp]
    \centering
    \includegraphics[width=\textwidth]{ranking_forest.pdf}
    \caption{The simplicity-preservation ranking of design levers.
    \emph{Panel A:} rows are levers sorted by pooled median
    $\rho = 1 - D(\text{lever})/D(\text{baseline})$; individual marks show
    each model $\times$ domain cell (color = model, shape = domain), with the
    pooled median and double-bootstrap $95\%$ CI per row; the shaded region
    $\rho < 0$ marks levers that backfire. Gold diamonds are human anchors
    computed from the moment-matched reconstructions
    (Section~\ref{sec:reconstruction}): AC $\rho = +0.62$, AC-B $+0.37$, menu
    $-0.07$; rungs without a diamond have no human experiment and are issued
    as lab predictions (Section~\ref{sec:humans}). Extreme negative
    matching-domain cells are clipped for display (see text). \emph{Panel B
    (inset):} cross-model rank-agreement (Kendall $\tau$) over the lever
    ordering. See Table~\ref{tab:ranking} for the underlying cells.}
    \label{fig:ranking}
\end{figure}

\paragraph{What the ranking localizes.}
Read from the bottom up, the assembled ranking localizes the constraint that
simple mechanism design relieves. Scaffolding strategizing---beliefs,
lookahead---makes play strictly worse, so the problem is not too little
strategic reasoning. Mechanically restating the rules is null on average and
harmful where it displaces the safety-relevant content, so the problem is not
raw comprehension bandwidth either. What works, wherever it appears and in
whatever syntactic form---the safety one-liners, the invariance-carrying menu
variant, the textbook strategy-proofness statement---is text that asserts the
invariant that makes truth-telling safe; and full OSP implementation, which
makes that invariant visible at every decision point, delivers the remainder.
As Section~\ref{sec:ranking-descriptions} put it: the binding cognitive
barrier in non-OSP mechanisms is not \emph{computing} the dominant strategy
given the game tree, but understanding the mechanism well enough to \emph{see
that truth-telling is safe}. The ordering, moreover, parallels the
theoretical hierarchy from which the taxonomy was derived ex ante
(Section~\ref{sec:framework}): the levers that theory says relax binding
cognitive constraints sit at the top, and the levers that theory says should
be irrelevant for strategy-proof mechanisms are, for agents who do not fully
trust that irrelevance, actively harmful. Section~\ref{sec:ranking-traces}
adds a final, mechanism-level piece of evidence: the effective descriptions
move bids \emph{without} changing what the agents say they are doing.
Section~\ref{sec:humans} then closes the loop, anchoring each rung to the
human experimental record where one exists and issuing the unanchored rungs
as predictions.

\subsection{What the Traces Say: Stated Reasoning and Realized Bids}
\label{sec:ranking-traces}

Every bid in our experiments is accompanied by a short free-text plan in
which the agent explains its choice (median 68 words). We analyze all
21{,}990 plans from the sealed-bid and clock cells with a transparent
dictionary of eighteen keyword features---stated intent to shade below or bid above
value, rehearsal of the second-price payment rule, dominance language,
worst-case and profit-margin rhetoric, and explicit mechanism confusion
(calling the auction ``first-price''), among others (full regular expressions
and per-cell prevalence in the replication package,
\path{analysis/build_trace_features.py}).\footnote{The plans are stated
strategies, not full chains of thought; all claims in this subsection concern
what models \emph{say}, with Table~\ref{tab:trace-dissociation} showing that
what they say tracks what they do.}

Two facts establish what the traces are. First, the modal plan applies
\emph{first-price} logic to the second-price auction: $53\%$ of sealed-bid
second-price traces state an intent to bid below value, $55\%$ worry about
``overpaying,'' and $31\%$ invoke a ``profit margin'' or ``buffer,'' while
only $2\%$ use dominance language and only $4\%$ correctly rehearse that the
winner pays the \emph{second}-highest bid; Claude explicitly mislabels the
mechanism ``first-price'' in $10\%$ of its traces, a mistake no other model
makes more than $0.1\%$ of the time. Second, the traces are not cheap talk:
among traces stating an intent to shade, $97.1\%$ of bids fall below value
(mean deviation $-\$1.84$); among traces stating an intent to bid above,
$87.6\%$ fall above ($+\$1.88$); and the largest errors come from the
4{,}541 traces that state no value-anchored intent at all (mean deviation
$-\$4.84$), whose plans anchor on uninformative quantities such as
``slightly above half my value.'' Stated direction almost always matches the
realized bid; the failure is in which plan is chosen. A pooled regression of
$|b - v|$ on the dictionary features nonetheless explains only a modest share
of the error variance (incremental $R^{2} \approx 0.075$ over model and
lever-family fixed effects); the regression table and robustness checks are
in Appendix~\ref{app:traces}.

Most important for interpreting the ranking, the levers of
Sections~\ref{sec:ranking-descriptions}--\ref{sec:ranking-scaffolds} produce
a complete \emph{dissociation} between stated reasoning and behavior
(Table~\ref{tab:trace-dissociation}): across every lever we test, language
and bids never move together. \emph{Payoff Safety} (B2) moves bids
substantially while leaving stated reasoning untouched: the invariant it
asserts (``your bid determines \emph{if} you win, not \emph{what} you pay'')
is echoed in zero of its 600 treated traces, and payment-rule rehearsal is
flat relative to baseline ($2.8\%$ treated vs.\ $3.9\%$, $p = 0.60$); the
plans keep the same shading rhetoric and simply shrink the shade to pennies
(``I will bid \$31.9 to avoid overbidding,'' value \$32). The
\emph{worst-case} scaffold (C1) does the opposite: it massively shifts
stated reasoning (worst-case language rises from $2\%$ to $43\%$) while bids
barely move ($3.24 \to 3.07$, $p = 0.81$). The \emph{Payoff Tree} is a
second bids-only lever: it cuts mean $|b-v|$ from $3.20$ to $1.29$ while
leaving payment-rule rehearsal exactly at its baseline rate ($33.0\%$ vs.\
$32.7\%$, $p = 0.90$)---an earlier reading on which the tree also moved
rule-rehearsal language was an artifact of the contaminated axis-2 baseline
cell documented in Section~\ref{sec:ranking-synthesis}, whose clock
narrative crowded out rule vocabulary. The \emph{menu} restatement moves
neither, serving as the placebo row. Because the first stage is null for
every verbalized-understanding measure, a classical mediation decomposition
bounds the share of Payoff Safety's behavioral effect that could run through
stated comprehension at roughly one third ($95\%$ upper bound $35\%$; point
estimate $\approx 0$), under assumptions the within-treatment association
(rule-rehearsing traces err \$1.5 less, CI $[-2.15,-0.56]$, non-causal)
cannot tighten. Stated comprehension and behavioral improvement are
separable outcomes, and our effective descriptions improve the latter
without the former---which is why we describe them as changing behavior, not
stated understanding, and why the traces should not be read as faithful
mediators of the lever effects.
\TODO{LLM-judge strategy taxonomy (trace option 1): validate the keyword
dictionary against a judged taxonomy on a 100-trace human-labeled set before
the ``2\% dominance recognition'' figure is made load-bearing.}

Finally, the traces speak to \emph{why} comprehension fails. In the GPT-4o
first-price and third-price variants of the same grid, the model deploys a
nearly identical shading script---shading language in $70\%/79\%/72\%$ of
second-/first-/third-price traces and mean bids of $0.88/0.87/0.84$ of
value---even though the optimum differs sharply across formats (truthful in
second-price; roughly $\tfrac{2}{3}v$ in our three-bidder first-price cell;
above value in third-price). The same words justify underbidding where
shading is optimal and where it is an error: the reasoning script is a
property of the model, not of the mechanism it faces.
\TODO{The cross-mechanism script-invariance result is currently
GPT-4o--only (recovered V12 logs); run FPSB/TPSB cells for Claude, Gemini,
and Gemma with the frozen dictionary before claiming generality.}

\begin{table}[htbp]
\centering
\small
\begin{tabular}{lcccc}
\toprule
& \multicolumn{2}{c}{Stated reasoning} & \multicolumn{2}{c}{Behavior} \\
\cmidrule(lr){2-3}\cmidrule(lr){4-5}
Lever (family) & target language & $\Delta$ vs.\ baseline &
mean $|b-v|$ & baseline \\
\midrule
Payoff Safety (B2) & payment-rule rehearsal & $3.9\%\to 2.8\%$ ($p=0.60$) &
$1.95$ ($p=0.015$) & $3.20$ \\
Worst-case scaffold (C1) & worst-case rhetoric & $2\%\to 43\%$ ($p=0.007$) &
$3.07$ ($p=0.81$) & $3.24$ \\
Payoff Tree (C1) & rule rehearsal & $33.0\%\to 32.7\%$ ($p=0.90$) &
$1.29$ ($p=0.017$) & $3.20$ \\
Menu restatement (B1) & --- & --- &
$3.55$ ($p=0.66$) & $3.20$ \\
\bottomrule
\end{tabular}
\caption{Dissociation between stated reasoning and bidding behavior
(pooled over four models, sealed-bid second-price cells; wild-cluster
bootstrap $p$-values clustered at the run level, with model-cluster
permutation cross-checks and duplicate-trace robustness in
\texttt{results/traces/mediation/}). No lever moves both columns: Payoff
Safety and the Payoff Tree improve bids without changing what agents say;
the worst-case scaffold changes what agents say without improving bids; the
menu restatement moves neither. The baseline is the corrected pooled axis
baseline (axis-1 and axis-3 cells only; the legacy axis-2 ``baseline'' cell
was itself a clock-framed description and is analyzed as a treatment in
Section~\ref{sec:ranking-descriptions}); under the dedicated-SPSB convention
the Payoff Safety contrast is the $-2.67 \to -1.53$ of
Section~\ref{sec:ranking-descriptions}. No language entry is reported for the
menu row because the menu prompt replaces the ``second-highest'' vocabulary
wholesale, so rule-rehearsal prevalence is zero by construction (prompt echo,
not comprehension).}
\label{tab:trace-dissociation}
\end{table}
